\documentclass[12pt]{article}
\usepackage[a4paper,margin=1in]{geometry}\usepackage[T1]{fontenc}
\usepackage{lmodern,microtype}\usepackage{amsmath,amssymb,amsthm,mathtools}
\usepackage{booktabs,array,enumitem,xcolor}\usepackage[numbers,sort&compress]{natbib}
\usepackage[colorlinks=true,linkcolor=blue!55!black,citecolor=blue!55!black,urlcolor=blue!55!black]{hyperref}
\linespread{1.07}
\newtheorem{theorem}{Theorem}[section]\newtheorem{proposition}[theorem]{Proposition}
\theoremstyle{remark}\newtheorem{remark}[theorem]{Remark}
\newcommand{\norm}[1]{\left\lVert#1\right\rVert}

\title{A Divisor-First Pivot Inside Gate A\\
\large Why Raw Projector Transport Is Sufficient but Not Necessary}
\author{Liang Wang}\date{July 2026}
\begin{document}\maketitle
\begin{abstract}
RH-202--RH-209 produce a split verdict.  The edge quartet has unique
conjugate branch labels and a coherent left/right quartic divisor, but naive
Haar transport fails for fixed and enlarged packets, and physical residues
do not share one scalar renormalization.  This paper determines the logical
route consequence.

We prove by an explicit rotating-similarity family that a spectral divisor
can remain exactly fixed while a selected Riesz projector moves by operator
norm one.  Thus convergence of raw projectors is a sufficient route to a
scalar determinant, but not a necessary condition for divisor stability.
An 81-angle audit keeps the characteristic coefficients fixed within
$2.84\times10^{-16}$ while the projector distance reaches one.

The physical evidence supports a divisor-first research order: preserve the
finite conjugate branch labels, densify and renormalize the quartic
coefficient flow, and keep the source-dependent residue cocycle in a
separate weighted ledger.  This is a pivot within Gate A, not its closure.
No coefficient limit, growing-cloud determinant, Fredholm realization,
Hilbert--P\'olya operator, or zeta-zero statement is obtained.
\end{abstract}

\section{The route decision forced by the batch}

The preceding layers establish four finite facts:
\begin{enumerate}[leftmargin=2em]
\item the naive Haar projector map has defect above one
\citep{WangRH202};
\item all four conjugate branch assignments are unique
\citep{WangRH204};
\item the left/right quartic coefficient discrepancy is below one percent
\citep{WangRH207};
\item increasing the top-modulus cloud to rank 32 does not repair the
two-sided angle \citep{WangRH209}.
\end{enumerate}

These facts do not say that the spectral direction is wrong.  They say that
one specific state-space realization is not the stable object visible in the
current data.

\section{A precise distinction}

Let $A_n$ be finite operators, $P_n$ selected Riesz projectors, and
\begin{equation}\label{eq:divisor}
 D_n(z)=\det(zI-A_n|_{\operatorname{Ran}P_n}).
\end{equation}
If transported restricted operators converge in trace norm, then their
finite determinants converge on compact sets under standard hypotheses
\citep{Simon2005}.  Hence strong state/operator transport is a sufficient
route.

The converse is false: determinants discard eigenvector geometry.  The next
theorem records this elementary but strategically important point.

\section{Rotating-projector counterexample}

Let
\begin{equation}\label{eq:rotation}
 U_\theta=\begin{pmatrix}\cos\theta&-\sin\theta\\
 \sin\theta&\cos\theta\end{pmatrix},\qquad
 D=\begin{pmatrix}0.7&0\\0&-0.2\end{pmatrix},
\end{equation}
and define
\begin{equation}\label{eq:family}
 A_\theta=U_\theta DU_\theta^*,\qquad
 P_\theta=U_\theta\begin{pmatrix}1&0\\0&0\end{pmatrix}U_\theta^*.
\end{equation}

\begin{theorem}[Divisor stability without projector stability]\label{thm:counterexample}
For every $\theta$,
\begin{equation}\label{eq:fixed-divisor}
 \det(zI-A_\theta)=(z-0.7)(z+0.2),
\end{equation}
while
\begin{equation}\label{eq:projector-drift}
 \norm{P_\theta-P_0}_2=|\sin\theta|.
\end{equation}
In particular the divisor is constant and the projector distance reaches one
at $\theta=\pi/2$.
\end{theorem}

\begin{proof}
Unitary similarity preserves the characteristic polynomial.  Direct
calculation of the two rank-one orthogonal projectors gives singular values
$|\sin\theta|$ for their difference.
\end{proof}

The example is normal and therefore does not exploit nonnormal instability.
It shows that the logical non-necessity is fundamental.

\section{Machine realization of the counterexample}

We sample 81 equally spaced angles from $0$ to $\pi/2$.  The maximum
coefficient-vector drift is $2.8306\times10^{-16}$, while the maximum
projector distance is exactly one to displayed precision.  This audit tests
the route logic and implementation; Theorem~\ref{thm:counterexample} itself
is exact.

\section{Physical evidence vector}

The route audit imports the following predeclared facts:
\begin{center}
\small
\begin{tabular}{p{0.58\textwidth}r}
\toprule
finite diagnostic & value\\
\midrule
largest naive Haar packet sine & $0.82388$\\
largest oblique-projector defect & $2.29068$\\
unique adjacent branch assignments & $4/4$\\
largest left/right quartic coefficient error & $0.008112$\\
largest common-scalar residue residual & $0.99985$\\
endpoint isolation feasibility cases & $6/6$\\
naive transport feasibility cases & $0/4$\\
expanded-cloud angle-gate cases & $0/24$ for $k>4$\\
\bottomrule
\end{tabular}
\end{center}
The scalar divisor is the only tested object combining basis invariance with
dual-channel coherence.  It is therefore the most economical next target.

\section{Three ledgers, not one}

We separate the program into:
\begin{description}[leftmargin=2.8em,style=nextline]
\item[Divisor ledger $Q$]
Eigenvalue multisets, monic coefficients, and unweighted Newton traces.
\item[State ledger $R$]
Right/left Riesz spaces, projectors, interlevel maps, and conditioning.
\item[Transfer ledger $W$]
Source--observation residues and weighted physical moments.
\end{description}
The current data support finite $Q$, diagnose difficulties in $R$, and show
that $W$ requires a branch-dependent cocycle.  None should be silently
substituted for another.

\section{The divisor-first route}

The revised order inside Gate A is:
\begin{enumerate}[leftmargin=2em]
\item retain the unique two-branch labels as the finite identity rule;
\item sample additional small-noise levels and both physical channels;
\item test intrinsic normalizations or recurrences for the full quartic
coefficient vector;
\item determine whether a growing family of divisors is normal and locally
uniform on a fixed complex domain;
\item only then seek a Fredholm/dynamical determinant realization;
\item return to state transport when such a realization specifies the right
renormalized coordinates.
\end{enumerate}

This is not permission to ignore operators permanently.  A dynamical
determinant ultimately needs trace-class, nuclear, symbolic, or equivalent
analytic structure.  The pivot only postpones an unsupported raw-projector
inductive limit.

\section{Why the fixed quartic is insufficient}

Even perfect convergence of one quartic would produce four limiting roots,
not a spectrum with a $T\log T$ count.  Gate A requires a growing cloud and
control of omitted factors.  The quartet is useful as:
\begin{enumerate}[leftmargin=2em]
\item a local branch-continuation seed;
\item a test of dual-channel universality;
\item a laboratory for coefficient renormalization;
\item a finite factor in a possible later determinant product.
\end{enumerate}
It is not the final determinant.

\section{What would falsify the pivot}

The divisor-first route should be abandoned or revised if denser levels show
any of:
\begin{enumerate}[leftmargin=2em]
\item loss of the branch labels or edge separation;
\item persistent left/right coefficient disagreement;
\item no compactness under any intrinsic predeclared normalization;
\item uncontrolled factor proliferation that destroys local uniformity;
\item dependence of the divisor on source/observation choices.
\end{enumerate}
Negative outcomes would still classify the finite quartet correctly as a
local diagnostic.

\section{Route coordinate and claim boundary}

The resulting coordinate is
\begin{equation*}
 \boxed{\texttt{finite\_dual\_channel\_divisor\_flow\_open\_renormalization}}.
\end{equation*}
Gate A is open.  Gates B--E remain untouched.  The paper proves only a
logical non-necessity theorem and records a finite evidence-based research
order.  It constructs no limiting divisor, Fredholm determinant,
self-adjoint operator, arithmetic trace identity, zeta divisor, or RH proof.

\bibliographystyle{plainnat}\bibliography{references}\end{document}
